\begin{theorem}[Six-contact estimate]\label{thm:height}
For every such $S,R$ and every $u\in\Q\setminus\Cset$,
\[
 h(u)\le H_6(S;R)M_S(u)^{24}.
\]
\end{theorem}
\begin{proof}
Take $B=(2,-3)$, the maximal order $\mathcal O=\Z[i,(1+j)/2]$, and
$\mu=2j+ij$. With its positive sign choice, the form
$\Trd(\mu^{-1}x\overline y)$ has matrix
\[
 \begin{pmatrix}0&0&-1&1\\0&0&1&-2\\1&-1&0&0\\-1&2&0&0\end{pmatrix}
\]
and determinant one, in the ordered basis
$(1,i,(1+j)/2,(i+ij)/2)$. The Riemann form is the convention of
\cite[equation (3) and Lemma 3]{linyang}; its Rosati involution is
$a\mapsto\mu^{-1}\overline a\mu$. This is the principal-polarization datum whose
normalizer quotient is $\mathbf P^1_j$, with branch values
$0,\infty,-16/27$ and indices $4,6,2$ \cite{babagranath,linyang}.
Put $K_0=\widehat{\mathcal O}^{\,*}$ and
$K_3=\ker(K_0\longrightarrow(\mathcal O/3\mathcal O)^*)$.
Let $N$ be the finite-adelic normalizer of $\widehat{\mathcal O}$.
It normalizes $K_3$. At a split prime the local normalizer is the
center times the maximal-order units; at each of the two division
primes its quotient by these groups has order two. Thus
$|N/(\mathbb A_f^*K_0)|=4$.
A central finite idele acts trivially on the level-three curve:
multiply it by a rational scalar to remove all finite valuations,
and choose the scalar's sign so that the residue at three is one.
The remaining idele lies in $K_3$, and the rational central scalar
acts trivially on the symmetric domain. Consequently the entire
level-three curve $Y_3$, with all its determinant components, admits
the action of $N/(\mathbb A_f^*K_3)$ with quotient $\mathbf P^1_j$,
and
\[
 \deg(Y_3\longrightarrow\mathbf P^1_j)
 \le4|(\mathcal O/3\mathcal O)^*|\le4\cdot81=324.
\]

We construct its integral model without changing this degree.
Apply \cite[Theorem 2.1 and Corollary 2.3]{buzzard} to $V_1(5)$
over $\mathbb Z[1/30]$. Its determinant image is full, its level
is maximal at the division primes, and $5\ge4$ is prime to six.
First add a full five-marking refining the $V_1(5)$-marking.
This gives the full-five moduli scheme, including every determinant
component, as a finite \'{e}tale cover. Next add the finite \'{e}tale sheaf
\[
 \operatorname{Isom}_{\mathcal O}
       ((\mathcal O/3\mathcal O),A[3]).
\]
This sheaf is nonempty on every geometric fiber. Indeed $B_3$ is
a division algebra and the rational Tate module is a rank-one
left $B_3$-module. Every $\mathcal O_3$-stable lattice in that module
is principal: choose a nonzero lattice element of least division
valuation and divide every other element by it on the right.
Thus the Tate lattice is free of rank one over $\mathcal O_3$.
One can also check this in characteristic zero and extend it to
all fibers by local constancy of the finite \'{e}tale $A[3]$ with its
$\mathcal O$-action.

The resulting smooth projective scheme $X_{35}$ carries full
markings at three and five. The group
$G_5=(\mathcal O/5\mathcal O)^*$ changes only the five-marking.
Its action is free: a stabilizer would give a polarized automorphism
preserving the full three-marking, hence a finite-order element of
$1+3M_4(\mathbb Z_3)$. That group is torsion-free, as follows by
expanding a hypothetical prime-order element at its first nonzero
$3$-adic term. The automorphism is therefore the identity, and
the full five-marking forces the group element to be the identity.
The quotient $X_{35}/G_5$ is smooth projective, its covering map is
finite \'{e}tale, and the universal abelian scheme, its canonical
principal polarization, its $\mathcal O$-action and its three-marking
descend. This quotient represents exactly the full-three moduli
problem, so its generic fiber is $Y_3$. The auxiliary five-marking
has contributed no degree factor to $Y_3\to\mathbf P^1_j$.

Each geometric component of $Y_3$ is a torsion-free cover of the
norm-one orbifold of Euler characteristic $-1/3$, and hence has genus
at least two. The normalizer action extends to its smooth proper
model by uniqueness of the stable model, including permutations of
components. Fix $p>324$ and work after strict henselization.
The action is tame. On each component use its effective action.
A finite-order automorphism of order prime to $p$ cannot specialize
to the identity, by considering its first nonzero term modulo a
power of the uniformizer. At a fixed point the same power-series
argument shows that the tangent character of the stabilizer is
faithful. Thus the stabilizer is cyclic.

For a nonidentity generator the tangent eigenvalue minus one is a
unit. Its fixed-point equation therefore has a unique lift by
Hensel's lemma. All powers of the generator fix that lift. It follows
that no new branch orbit appears and that distinct branch orbits
cannot collide: their stabilizers at a purported common specialization
would belong to one cyclic group, whose nontrivial elements have the
same unique fixed lift. Averaging a parameter gives a local invariant
ring $R[[z^e]]$. Hence the quotient is smooth and retains the three
disjoint branch sections and their indices $4,6,2$.

Its genus-zero model with these three ordered sections is isomorphic
to $\mathbf P^1_{\mathbb Z_p}$ carrying the ordered triple
$(0,\infty,-16/27)$, which is disjoint since $2$ and $3$ are units.
The isomorphism agrees with $j$ on the generic fiber, by uniqueness
of a projective coordinate on three points. Thus it is the integral
model in the stated $j$-coordinate; no further denominator prime
has been discarded.

Pull back by $J(u)=-57600F(u)^2/G(u)^6$. If
\[
 H(u)=(u^2-60u-1800)(13u^2+480u+3600)(11u^2+600u+7200),
\]
direct calculation gives $G^6-97200F^2=H^2$. The pullback indices are
$2,6,2$. Tame normalization replaces $e$ by $e/\gcd(e,a)$ and leaves
$2,1,1$. The discriminants and resultants have prime support at most 13,
so this persists at $p>324$. The normalized cover of $\mathbf P^1_u$
has degree at most 324, not $12\cdot324$, and is branched only at the
six contacts. Locally it has equation $u-c=z^2$ after unramified
extension. Positive even contact valuation is unramified; negative
valuation specializes to the unbranched point at infinity. This is
exactly the support specified by $M_S(u)$.

Let $K$ be the splitting field of the fiber over $u$ of the normalized
cover. The fiber is finite \'{e}tale of degree at most $324$, so $K/\mathbb Q$
is normal with $[K:\mathbb Q]\le F_0=324!$. The preceding local
calculation makes it unramified outside
$T=S\cup\operatorname{supp}M_S(u)$. Choose a fiber $A/K$.
The proper model and its universal abelian scheme give good reduction
outside the places over $S$.

All conjugates of $A$ have the same geometric polarized moduli point.
Indeed Baba--Granath's polarized equivalence relation gives
$V_{\mathbb Z}=V/W$; for discriminant six its two twisting involutions
generate the whole Atkin--Lehner group $W$. On the finite nonzero
$j$-locus the polarized point is equivalently specified by
\[
 [J_2,J_4,J_6,J_{10}]
 =[12(j+1),6(j^2+j+1),4(j^3-2j^2+1),j^3].
\]
Here $j=J(u)$ is rational, finite and nonzero: $F(u)\ne0$ and
$\operatorname{disc}(G)=-10800<0$.
